\usepackage{mathrsfs}
\let\mathcal\mathscr
\multlinegap0pt
\newcommand{\Psfrac}[2]{(\sfrac{#1}{#2})}
\newcommand{\Ppsfrac}[2]{(\sfrac{(#1)}{#2})}
\newcommand{\psfrac}[2]{\sfrac{(#1)}{#2}}
\newcommand{\spfrac}[2]{\sfrac{#1}{(#2)}}
\newcommand{\pspfrac}[2]{\sfrac{(#1)}{(#2)}}

\let\epsilon\varepsilon
\newcommand{\quand}{\quad\text{and}\quad}
\DeclareMathOperator{\Ai}{Ai}
\newcommand{\Changel}{\mathcode`l="7160}
\newcommand{\Changelback}{\mathcode`l="716C}
\newcommand{\NoChangel}[1]{%
\expandafter\let\csname old\string#1\endcsname=#1
\let#1=\relax
\newcommand{#1}{\mathcode`l="716C\csname old\string#1\endcsname\mathcode`l="7160 }%
}

\NoChangel{\log}
\NoChangel{\ln}
\NoChangel{\lim}
\NoChangel{\limsup}
\NoChangel{\liminf}
\NoChangel{\varlimsup}
\NoChangel{\varliminf}
\NoChangel{\varinjlim}
\NoChangel{\varprojlim}



\usepackage{mathtools}

\newcommand\N{{\mathbb N}}
\newcommand\R{{\mathbb R}}
\newcommand\Z{{\mathbb Z}}

\theoremstyle{plain}
\newtheorem{thm}{Theorem}
\newtheorem{lemma}{Lemma}
\newtheorem{prop}{Proposition}

\theoremstyle{definition}
\newtheorem{rmq}{Remark}[section]

\newcommand\e{{\epsilon}}
\newcommand{\DomainXYT}{\mathcal{R}}

